\documentclass[11pt,a4paper]{article}
\usepackage[T1]{fontenc}
\usepackage[utf8]{inputenc}
\usepackage{lmodern,amsmath,amssymb}
\usepackage[margin=25mm]{geometry}
\usepackage[hidelinks]{hyperref}
\hypersetup{pdftitle={A Lorentz-invariant background and the state-route constant: a conditional link to the rad},pdfauthor={navstokgap research working notes}}
\newcommand{\tightlist}{\setlength{\itemsep}{0pt}\setlength{\parskip}{0pt}}
\setlength{\emergencystretch}{3em}
\setcounter{secnumdepth}{0}
\begin{document}
\hypertarget{a-lorentz-invariant-background-and-the-state-route-constant-a-conditional-link-to-the-radiation-unit}{%
\section{A Lorentz-invariant background and the state-route constant: a
conditional link to the radiation
unit}\label{a-lorentz-invariant-background-and-the-state-route-constant-a-conditional-link-to-the-radiation-unit}}

\textbf{Result, 2026-09-27; revised the same day after an adversarial
Fable review.} STATE item 2 asked for a physical reason why the constant
\(\zeta\) of the state route (GPT-6 Astra's
\href{newton-indeterminacy-routes.md}{routes note}, floor with
\(h_*=2\zeta\); the \href{principia-fifth-postulate.md}{fifth-postulate
note}, Theorem B\('\)) should be a fixed multiple of the radiation unit
\(h_{\rm rad}\) of Theorem U
(\href{necessity-unit-and-indeterminacy.md}{unit-and-indeterminacy
note}). Stochastic electrodynamics supplies one, conditionally.

\begin{enumerate}
\def\labelenumi{\arabic{enumi}.}
\tightlist
\item
  A stationary random electromagnetic background that is Lorentz
  invariant has energy \(\kappa\omega\) per normal mode, one constant
  \(\kappa\) with the units of action.
\item
  For a charged oscillator in equilibrium with it, the \emph{resonant}
  part of the motion is Gaussian with \(\sqrt{\det\Sigma}=\kappa\),
  independent of mass, charge and frequency (Theorem 1). The full
  momentum variance diverges in the Lorentz-invariant background; a
  frequency cutoff removes the divergence and breaks the invariance that
  fixed the spectrum. So the background realizes the restriction of
  Theorem B\('\) with \(\zeta=\kappa\) for band-filtered (resonant)
  variables.
\item
  \emph{If} the Planck spectrum, zero-point term included, is derived
  within stochastic electrodynamics, the thermal quantum is twice the
  zero-point constant: \(h_P=4\pi\kappa\), where \(h_P\) is Planck's
  constant. With \(2\zeta=2\kappa\) this gives \(2\zeta=h_P/(2\pi)\).
  Relative to Theorem U's \(h_{\rm rad}\), which differs from \(h_P\) by
  the pure number \(h_P=(\pi^4/15)^{1/3}h_{\rm rad}\) stated in the unit
  note, \[2\zeta=\gamma\,h_{\rm rad},\qquad
  \gamma=\frac{(\pi^4/15)^{1/3}}{2\pi}\approx0.297 .\]
\end{enumerate}

Every stochastic-electrodynamics derivation of the Planck spectrum is
disputed in that literature (§3), and the value \(\kappa=\hbar/2\) is
then experimental, fixed by matching the zero-point term. The link is
therefore conditional on step 3 and restricted to resonant variables by
step 2. Astra's closure premise, that records cannot read the background
and leave its equilibrium set, stays open. The ingredients are
established; the assembly is this note's, with no novelty claimed for
the parts.

\hypertarget{the-background}{%
\subsection{1. The background}\label{the-background}}

A homogeneous, isotropic, stationary random electromagnetic field whose
spectrum is Lorentz invariant has energy density proportional to
\(\omega^3\), that is, energy per normal mode linear in frequency,
\(\kappa\omega\). Boyer states it in the abstract of his derivation
(\href{https://doi.org/10.1103/PhysRev.182.1374}{Boyer 1969}, abstract:
the zero-point spectrum is ``linear in frequency'',
\(\frac12\hbar\omega\) per normal mode); the random-electrodynamics
programme goes back to
\href{https://doi.org/10.1098/rspa.1963.0220}{Marshall (1963)}
(metadata). Lorentz invariance leaves \(\kappa\) free.

\hypertarget{theorem-1-the-resonant-gaussian-area}{%
\subsection{2. Theorem 1: the resonant Gaussian
area}\label{theorem-1-the-resonant-gaussian-area}}

\textbf{Theorem 1.} Let a particle of mass \(M\) and charge \(e\) be
bound harmonically at frequency \(\omega_0\) with radiation damping,
\(M\ddot x=-M\omega_0^2x+M\tau\dddot x+eE_x\), \(\tau=2e^2/(3Mc^3)\),
weak damping \(\tau\omega_0\ll1\), driven by the Gaussian background of
§1. The stationary state is Gaussian, and its resonant part (frequencies
within a few linewidths of \(\omega_0\)) has

\[\langle x^2\rangle_{\rm res}=\frac\kappa{M\omega_0},\qquad
\langle p^2\rangle_{\rm res}=\kappa M\omega_0,\qquad
\langle xp+px\rangle=0,\qquad\sqrt{\det\Sigma_{\rm res}}=\kappa,\]

independent of \(M\), \(e\) and \(\omega_0\).

\emph{Proof.} Linear response to Gaussian forcing is Gaussian. Near
resonance the susceptibility satisfies
\(|\chi|^{-2}\simeq4\omega_0^2(\omega-\omega_0)^2+\tau^2\omega_0^6\),
whose integral gives \(\int|\chi|^2d\omega\simeq\pi/(2\tau\omega_0^4)\);
with the background's spectral density the charge and \(\tau\) cancel,
and the resonant mean energy is the field energy per mode at
\(\omega_0\), \(\kappa\omega_0\) (Planck's resonator relation for a
dipole with radiation damping,
\href{https://doi.org/10.1002/andp.19003060105}{Planck 1900}, metadata),
split equally between kinetic and potential terms. Stationarity gives
\(\langle xp+px\rangle=0\). (Hand check of the resonant moments by the
reviewer.) \(\square\)

\textbf{The off-resonant part.} In the \(\omega^3\) background the
momentum variance has a contribution from \(\omega_0\ll\omega\ll1/\tau\)
of relative size about \(1/(\pi\tau\omega_0)\), which grows as the
damping weakens, and above \(1/\tau\) the integrand behaves like
\(1/(\tau^2\omega)\): the full \(\langle p^2\rangle\) diverges
logarithmically. The position moment is safe; the off-resonant part of
\(\langle x^2\rangle\) is of relative order
\(\tau\omega_0\ln(1/\tau\omega_0)\). This is the known difficulty of the
oscillator in stochastic electrodynamics
(\href{https://doi.org/10.1007/BF00728140}{Goedecke 1983}, metadata;
\href{https://doi.org/10.3390/atoms7020059}{Nieuwenhuizen 2019},
abstract: only ``after introducing a cut-off in the stochastic power
spectrum and regularizing the stochastic force'' are the integrals
dominated by resonance). A cutoff \(\Omega_c\ll\sqrt{\pi\omega_0/\tau}\)
keeps the resonant answer, and any such cutoff breaks the Lorentz
invariance of §1, as Astra's routes note already observes. Theorem 1 is
therefore a statement about band-filtered variables.

\textbf{The link to Theorem B\('\).} The Gaussian states with
\(\sqrt{\det\Sigma}\ge\kappa\) form an affine-symplectic-invariant,
noise-closed set, and Theorem 1 puts the resonant variables of every
harmonically bound body on its lower boundary. The noise direction that
Theorem B\('\) needs is supplied by a forgotten record, which adds
momentum variance only (\href{newton-record-parallel-move.md}{record
note}, Proposition 1).

\hypertarget{the-radiation-unit-conditionally}{%
\subsection{3. The radiation unit,
conditionally}\label{the-radiation-unit-conditionally}}

Boyer derives the Planck law ``without the formalism of quantum theory''
from this background, classical electrodynamics of dipole oscillators
and classical equipartition of the particles' kinetic energy (Boyer
1969, abstract), by an Einstein--Hopf argument. What such a derivation
fixes is a ratio: the thermal part has the form of Planck's with quantum
\(h_P\), and the zero-point part is \(\kappa\omega=h_P\omega/(4\pi)\),
so \(h_P=4\pi\kappa\); the value \(\kappa=\hbar/2\) is then read from
experiment.

The derivation is contested. Senatchin argues that ``his derivation
contains a loophole in its argument'', the wall damping making the
equilibrium radiation inhomogeneous
(\href{https://arxiv.org/abs/physics/0105054}{arXiv:physics/0105054},
passage as reported by the reviewer), citing
\href{https://doi.org/10.1119/1.12221}{Jiménez, de la Peña and Brody
(1980)} (metadata); Boyer's own historical review maintains the
calculation with a modification of the wall damping
(\href{https://arxiv.org/abs/1711.04179}{arXiv:1711.04179}, passage as
reported). Nonlinear systems in stochastic electrodynamics are known to
depart from the Planck spectrum (reviewer's recollection, not checked
here). So the honest statement is: within stochastic electrodynamics,
and conditional on an SED derivation of the Planck spectrum with its
zero-point term,

\[2\zeta=\frac{h_P}{2\pi}=\frac{(\pi^4/15)^{1/3}}{2\pi}\,h_{\rm rad}\approx0.297\,h_{\rm rad}.\]

\hypertarget{consequence-for-state}{%
\subsection{4. Consequence for STATE}\label{consequence-for-state}}

STATE item 2 asked for the identification \(2\zeta=\gamma h_{\rm rad}\)
and for closure under recording. The identification holds within
stochastic electrodynamics, conditionally on its disputed derivation of
the Planck spectrum and for resonant variables, with
\(\gamma=(\pi^4/15)^{1/3}/(2\pi)\). Closure under recording remains the
open premise, now joined by the tension between a Lorentz-invariant
background and the cutoff that the oscillator's momentum variance
requires.
\end{document}
